% ============================================================
%  rvi_analysis.tex
%  Technical note: RVI policy heatmaps for a 2-server / 2-type
%  queue under queue-lateness and binary deadline rewards.
% ============================================================
\documentclass[11pt,a4paper]{article}

\usepackage{amsmath,amssymb,amsthm}
\usepackage{geometry}
\geometry{margin=2.5cm}
\usepackage{booktabs}
\usepackage{array}
\usepackage{caption}
\usepackage{hyperref}
\usepackage{xcolor}
\usepackage{microtype}
\usepackage{parskip}
\usepackage{verbatim}

\newtheorem{remark}{Remark}
\newtheorem{observation}{Observation}

\title{Relative Value Iteration for a Two-Server Scheduling Problem:\\
       Policy Structure, Heatmap Enumeration, and the Role\\
       of the Reward Function}
\author{Technical Note}
\date{}

\begin{document}
\maketitle

\begin{abstract}
We study the optimal scheduling policy for a two-server, two-job-type fully
flexible queueing system with asymmetric holding costs and a fixed physical
deadline, solved via Relative Value Iteration (RVI).
We introduce an \emph{enumeration-based} heatmap that directly constructs
canonical states and queries the policy map without simulation, eliminating
all sampling artefacts.
Under the queue-lateness (QL) reward the heatmap reveals a clean three-region
threshold structure whose boundary, in physical time, converges to the deadline
$T_\text{real}$ as the tick rate increases.
We verify the Bellman average-cost value $g^*$ against an independent simulation
estimate and demonstrate that both agree within Monte Carlo error; we also clarify
the relationship between the per-RVI-step cost and the per-real-event cost used
by the policy comparer.
We then analyse the relationship between the discrete uniformised chain and the
underlying continuous-time system, showing that tick events create no additional
scheduling decisions and that the tick rate acts as a FIL-resolution parameter:
the optimal threshold converges to the physical deadline $T_\text{real}$ as the
tick rate increases.
Finally, we give an elaborate explanation of why the binary deadline reward
produces noisy heatmaps: the cost function's near-zero gradient away from the
deadline creates action-value near-ties across most of the state space, and the
\texttt{g\_stable} convergence fallback does not provide sufficient resolution
to break these ties deterministically.
\end{abstract}

\tableofcontents
\newpage

% ============================================================
\section{System Model}
% ============================================================

Consider a queueing system with the following structure.

\paragraph{Servers.}
There are $K = 2$ server pools, each containing a single server.  Both
pools are fully flexible: each server can process either job type at the same
exponential service rate $\mu$.  In the numerical study we set $\mu = 0.4$.

\paragraph{Job types.}
There are $N = 2$ job types.  Jobs of type $n$ arrive as an independent
Poisson stream with rate $\lambda_n$.  We set $\lambda_0 = \lambda_1 = 0.1$.
Upon arrival, a job waits in its own queue until it is assigned to an idle
server.

\paragraph{Waiting times (FIL).}
The system tracks, for each type $n$, the waiting time of the
\emph{first-in-line} (FIL) job.  A discrete tick event (rate $\gamma$) advances
all waiting times by one unit simultaneously.  The physical time per tick is
$1/\gamma$, and the physical deadline is $T_\text{real} = 3$ time units, giving
a due-tick count of $d = T_\text{real} \cdot \gamma$.

\paragraph{Traffic intensity.}
\[
  \rho \;=\; \frac{\lambda_0 + \lambda_1}{K\mu}
          \;=\; \frac{0.1 + 0.1}{2 \times 0.4}
          \;=\; 0.25.
\]

\paragraph{Cost rates.}
Asymmetric holding cost rates $c_0 = 3$ and $c_1 = 1$ are applied, so
type-0 jobs are three times more expensive to delay.

\paragraph{Queue-depth formula.}
Using the geometric tail bound $\rho^{k+1} < \varepsilon$ with
$\varepsilon = 0.01$ and $\rho = 0.25$:
\[
  k_\text{true}
    = \left\lceil \frac{\log \varepsilon}{\log \rho} \right\rceil
    = \left\lceil \frac{\log 0.01}{\log 0.25} \right\rceil
    = \lceil 3.322 \rceil = 4.
\]
Thus a queue depth of 4 positions per job type captures 99\% of the
probability mass; the FIL-only projection ($k=1$) is evaluated against this
richer model in Section~\ref{sec:results}.


% ============================================================
\section{MDP Formulation}
\label{sec:mdp}
% ============================================================

\subsection{Uniformised Markov Decision Process}

The continuous-time dynamics are embedded into a uniformised discrete-time MDP
via the standard uniformisation trick.  At each step exactly one of the
following events occurs:
\begin{itemize}
  \item \textbf{Arrival} of a job of type $n$ (probability $\propto \lambda_n$);
  \item \textbf{Tick} --- all waiting times advance by one unit
        (probability $\propto \gamma$);
  \item \textbf{Service completion} of a job of type $n$ at server pool $k$
        (probability $\propto \mu \cdot b_{k,j(n)}$, where $b_{k,j}$ is the
        current busy count);
  \item \textbf{Nothing} --- a self-loop that keeps the total event rate
        constant.
\end{itemize}
This yields a uniformisation rate
$\Lambda = \sum_n \lambda_n + \gamma + \mu \sum_{k,j} b_{k,j}$.

\subsection{State Space}

A state $s \in \mathcal{S}$ is a tuple
\[
  s = \bigl(f_0,\, f_1,\; B,\; \alpha,\; \kappa \bigr),
\]
where
\begin{itemize}
  \item $f_n \in \{-1, 0, 1, \ldots\}$ is the FIL waiting time of type $n$
        ($f_n = -1$ means no job is waiting);
  \item $B = (b_{k,j})$ with $b_{k,j} \in \{0,1\}$ is the busy configuration
        ($b_{k,j} = 1$ iff server pool $k$ is currently serving a job of
        type $\text{can\_serve}[k][j]$);
  \item $\alpha \in \{0, 1, \ldots\}$ is the \emph{action counter} --- the
        current position in the ordered action queue;
  \item $\kappa \in \{\textsc{AwaitEvent},\, \textsc{AwaitAction}\}$ is the
        state category.
\end{itemize}

\subsection{Action Queue and Action Counter}

When at least one server is idle and at least one job is waiting, a
sequence of assignment decisions must be made before the next event is drawn.
The candidate assignments are collected into an ordered \emph{action queue}
sorted by descending FIL (FIFO order), with ties broken by type index.

At category $\kappa = \textsc{AwaitAction}$ and counter $\alpha$, the
decision-maker faces the action at position $\alpha$ in the queue:
\begin{itemize}
  \item \textbf{Action 1 (assign):} assign the job at position $\alpha$ to the
        corresponding idle server; remove infeasible actions from the queue;
        re-evaluate the counter.
  \item \textbf{Action 0 (skip):} pass on position $\alpha$ and advance to
        $\alpha + 1$.  If $\alpha$ reaches the end of the queue, no further
        assignments are made and the system transitions to \textsc{AwaitEvent}.
\end{itemize}

\begin{remark}
\textbf{Action 0 is not idling.}  Choosing to skip the FIFO-top candidate
at counter $\alpha = 0$ simply means the policy prefers to assign a
lower-ranked job (e.g.\ type 0 instead of type 1) when the two FIL values
are in the ``wrong'' order for the cost structure.
\end{remark}

\subsection{Immediate Cost}

Cost is incurred only at \textsc{AwaitEvent} states (no physical time passes
during action steps).  Two reward types are studied:

\paragraph{Queue-lateness (QL) reward (reward\_type = 1):}
\begin{equation}
  c_\text{QL}(s)
    = \sum_{n=0}^{1} c_n \cdot \max\!\bigl(0,\; f_n - d_n\bigr).
  \label{eq:ql_cost}
\end{equation}
The cost is zero until the deadline and grows linearly with overdue-ness
thereafter.

\paragraph{Binary deadline reward (reward\_type = 0):}
\begin{equation}
  c_\text{bin}(s)
    = \sum_{n=0}^{1} c_n \cdot \mathbf{1}\!\bigl\{f_n > d_n\bigr\}.
  \label{eq:bin_cost}
\end{equation}
The cost is a flat penalty that fires as soon as a job exceeds its deadline,
with no further increase for additional overdue-ness.


% ============================================================
\section{Relative Value Iteration}
\label{sec:rvi}
% ============================================================

\subsection{Average-Cost Bellman Equations}

We seek the optimal average cost per step $g^*$ and the relative value
function $h : \mathcal{S} \to \mathbb{R}$ satisfying the Bellman equations:
\begin{align}
  h(s) &= c(s) - g^* + \sum_{s'} P(s' \mid s)\, h(s'),
         \quad s \in \mathcal{S}_\text{event},
  \label{eq:bell_event} \\[4pt]
  h(s) &= \min_{a \in \{0,1\}} \sum_{s'} P(s' \mid s, a)\, h(s'),
         \quad s \in \mathcal{S}_\text{action}.
  \label{eq:bell_action}
\end{align}
The function $h$ is unique up to an additive constant; we normalise by fixing
$h(s_\text{ref}) = 0$ for a reference state $s_\text{ref}$.

\subsection{State-Space Truncation}

The FIL values are clamped to a finite truncation level $M$: any state with
$f_n > M$ is treated as $f_n = M$ (a boundary self-loop).  All reachable states
with $f_n \leq M$ are enumerated by breadth-first search (BFS) from the initial
empty-queue state.

An automatic procedure selects $M$: RVI is solved at an initial guess,
$M$ is incremented by 2, and RVI is re-solved until the relative change in
$g^*$ falls below a tolerance (here 0.01).

\subsection{The RVI Update}

Starting from $h^{(0)} \equiv 0$, iterate:
\begin{align*}
  h^{(t+1)}(s) &=
  \begin{cases}
    c(s) + \displaystyle\sum_{s'} P(s' \mid s)\, h^{(t)}(s'),
      & s \in \mathcal{S}_\text{event}, \\[6pt]
    \displaystyle\min_{a \in \{0,1\}}
      \sum_{s'} P(s' \mid s, a)\, h^{(t)}(s'),
      & s \in \mathcal{S}_\text{action},
  \end{cases} \\[4pt]
  g^{(t+1)} &= h^{(t+1)}(s_\text{ref}), \\
  h^{(t+1)}(s) &\leftarrow h^{(t+1)}(s) - g^{(t+1)}.
\end{align*}

\subsection{Convergence Criterion}

The theoretically correct criterion is the \emph{span seminorm}:
\[
  \text{span}\!\left(h^{(t+1)} - h^{(t)}\right)
    = \max_s \Delta^{(t)}(s) - \min_s \Delta^{(t)}(s)
    \xrightarrow{t\to\infty} 0,
\]
where $\Delta^{(t)}(s) = h^{(t+1)}(s) - h^{(t)}(s)$.  Unlike
$\max_s |\Delta^{(t)}(s)|$, the span is not inflated by a near-constant offset
added by boundary self-loops at $f_n = M$.

In practice, the truncation at $M$ prevents the span from reaching exactly
zero.  A fallback \emph{g-stable} criterion fires when
$\lvert g^{(t+1)} - g^{(t)} \rvert < 10^{-10}$ for five consecutive
iterations; $g^*$ converges reliably even when the span does not.

\subsection{Action Extraction}

After convergence, the optimal action for every \textsc{AwaitAction} state
is stored:
\begin{equation}
  \pi^*(s)
    = \operatorname*{arg\,min}_{a \in \{0,1\}}
      Q(s, a),
  \qquad
  Q(s, a) = \sum_{s'} P(s' \mid s, a)\, h^*(s').
  \label{eq:argmin}
\end{equation}
The function $h^*$ is then \emph{discarded}; the only persistent object
is the \emph{action map}
\[
  \texttt{action\_map} : \text{encode}(s) \;\mapsto\; \pi^*(s) \in \{0,1\}.
\]


% ============================================================
\section{State Encoding and the Action Map}
\label{sec:encoding}
% ============================================================

\subsection{The Encoder}

A single 64-bit integer $k(s)$ is computed via mixed-radix encoding:
\[
  k(s)
  = \underbrace{\tilde{f}_0}_{\text{FIL}_0}
    + (M{+}2)\,\underbrace{\tilde{f}_1}_{\text{FIL}_1}
    + \ldots
    + \prod_n (M{+}2)
      \cdot \biggl[
        \sum_{k,j} b_{k,j} \cdot \sigma_{k,j}
        + \sigma_\kappa \cdot \kappa
        + \sigma_\alpha \cdot \alpha
      \biggr],
\]
where $\tilde{f}_n = \min(f_n, M) + 1$ (so that $f_n = -1$ maps to 0) and
$\sigma_{\cdot}$ are stride constants determined by the range of each component.

\paragraph{Critical observation.}
The busy configuration $B$ and the action counter $\alpha$ are both part of
the key.  Two states that share the same FIL values but differ in $B$ or
$\alpha$ receive \emph{different} keys and have \emph{independently stored}
optimal actions.  In particular:
\begin{align*}
  k\!\bigl(f_0, f_1,\; b_{0,0}{=}1,\; \alpha{=}0\bigr)
  &\neq
  k\!\bigl(f_0, f_1,\; b_{0,1}{=}1,\; \alpha{=}0\bigr),
\end{align*}
i.e.\ whether pool 0 is currently serving a type-0 or a type-1 job
leads to different entries in the map and potentially different optimal
decisions.

\subsection{What the Action Map Stores}

The action map does \emph{not} store $h^*(s)$.  It stores exactly one bit
per state: the argmin action $\pi^*(s) \in \{0, 1\}$.  The $h$-values
only influence the map through \eqref{eq:argmin} at construction time.

For a 2-server fully flexible system with $M = 15$ the map contains
approximately 3\,264 entries corresponding to all \textsc{AwaitAction} states
reached by the BFS.


% ============================================================
\section{Enumerated Policy Heatmap}
\label{sec:heatmap}
% ============================================================

\subsection{Motivation: Limitations of Simulation-Based Heatmaps}

A simulation-based heatmap runs the MDP under a given policy for a large number
of steps, records the action taken at each state in a $(f_0, f_1)$ grid, and
marks unvisited cells with ``\texttt{-}''.  This approach has two failure modes:

\begin{enumerate}
  \item \textbf{Unvisited cells (\texttt{-}).}
    Rare states (high FIL, unusual busy configuration) are never visited in a
    finite simulation, leaving gaps in the map.

  \item \textbf{Sampling-averaged inconsistency.}
    The same $(f_0, f_1)$ cell may be visited under different busy
    configurations $B$ and action counters $\alpha$.  The last recorded action
    overwrites earlier ones, so the cell reflects whichever configuration was
    sampled most recently --- a simulation artefact, not the policy.
\end{enumerate}

\subsection{The Enumeration Approach}

For each $(f_0, f_1) \in \{0,\ldots,M_\text{map}\}^2$, construct the
\emph{canonical} state
\begin{equation}
  s^\dagger(f_0, f_1)
  = \bigl(f_0,\; f_1,\; b_{0,0} = 1 \text{ (all others 0)},\;
    \alpha = 0,\; \kappa = \textsc{AwaitAction}\bigr).
  \label{eq:canonical}
\end{equation}
This represents: pool 0 busy on a type-0 job, pool 1 idle, both job types
waiting, first assignment decision.

\paragraph{Procedure for each cell $(f_0, f_1)$:}
\begin{enumerate}
  \item Set FIL values: $f_n$ entered directly into the queue manager.
  \item Set busy configuration: $b_{0,0} = 1$, all others 0.
  \item Call \texttt{generate\_actions($f_0, f_1$)}: fills the action queue
        $\mathcal{A}$ sorted by descending FIL (FIFO order), ties broken by
        type index.
  \item Set $\alpha = 0$.
  \item Compute key $k = \text{encode}(s^\dagger)$ and look up
        $a = \texttt{action\_map}[k]$ (fallback: $a = 1$ if key absent).
  \item \textbf{Display:}
    \begin{itemize}
      \item $a = 0$: show ``\texttt{.}'' (skip FIFO-top candidate).
      \item $a = 1$: show $\mathcal{A}[0].\text{job\_type}$
            (type of the FIFO-top candidate).
    \end{itemize}
\end{enumerate}

\paragraph{Result.}  Every cell $(f_0, f_1)$ has a definite, artefact-free
value.  There are no ``\texttt{-}'' cells.

\subsection{The Action-Queue at the Canonical State}

With pool 1 idle and both types waiting, \texttt{generate\_actions} yields:
\begin{equation}
  \mathcal{A} =
  \begin{cases}
    \bigl[(\text{pool 1, type 1}),\; (\text{pool 1, type 0})\bigr]
      & \text{if } f_1 > f_0, \\[3pt]
    \bigl[(\text{pool 1, type 0}),\; (\text{pool 1, type 1})\bigr]
      & \text{if } f_0 \geq f_1.
  \end{cases}
  \label{eq:action_queue}
\end{equation}

Consequently, ``\texttt{.}'' above the diagonal ($f_1 > f_0$) means:
\emph{RVI skips type-1 even though it is older; at counter $\alpha=1$ the
policy will assign pool 1 to type-0 instead.}  This is the cost-priority
override.

\subsection{Limitation: One Canonical Busy Configuration}

The canonical state \eqref{eq:canonical} fixes pool 0 busy on type 0.  The
action map contains separate entries for every other busy configuration.
In particular, the configuration ``pool 0 busy on type 1'' (key:
$b_{0,1} = 1$) may yield different optimal actions for the same $(f_0, f_1)$
because type 1 is already being handled, reducing the urgency of assigning
pool 1 to type 1 from pool 1.

For the fully symmetric system studied here (identical pools, equal $\lambda$
and $\mu$) the canonical state is representative: swapping pool labels is a
symmetry of the problem.  The qualitative threshold structure is expected to
be the same across all symmetric busy configurations, with the threshold
potentially shifting by $O(1)$ tick in asymmetric configurations.


% ============================================================
\section{Results: Queue-Lateness Reward}
\label{sec:results}
% ============================================================

\subsection{Three-Region Policy Structure}

\begin{table}[h]
\centering
\caption{Three regions of the RVI policy heatmap (QL reward, tick\_rate = 4).}
\label{tab:regions}
\begin{tabular}{lcl}
\toprule
Region & Cell symbol & Interpretation \\
\midrule
$f_0 \geq f_1$ (on/below diagonal) & \texttt{0}
  & Type 0 dominates in both age \emph{and} cost; always served. \\
Above diagonal, small excess & \texttt{.}
  & Type 1 is older but cost advantage of type 0 wins; skip type 1. \\
Above diagonal, large excess & \texttt{1}
  & Type 1 is so much older that its urgency overcomes the 3$\times$ cost gap. \\
\bottomrule
\end{tabular}
\end{table}

The resulting pattern (tick\_rate $= 4$, due $= 12$) is shown below in
condensed form (columns 0--19):

\begin{verbatim}
      FIL_1:  0  1  2  3  4  5  6  7  8  9 10 11 12 13 14 15 16 17 18 19
FIL_0
0     :        0  .  .  .  .  .  .  .  .  .  .  .  .  1  1  1  1  1  1  1
1     :        0  0  .  .  .  .  .  .  .  .  .  .  .  .  1  1  1  1  1  1
2     :        0  0  0  .  .  .  .  .  .  .  .  .  .  .  1  1  1  1  1  1
3     :        0  0  0  0  .  .  .  .  .  .  .  .  .  .  .  .  1  1  1  1
4     :        0  0  0  0  0  .  .  .  .  .  .  .  .  .  .  .  .  1  1  1
5     :        0  0  0  0  0  0  .  .  .  .  .  .  .  .  .  .  .  .  .  1
6     :        0  0  0  0  0  0  0  .  .  .  .  .  .  .  .  .  .  .  .  .
  ...
\end{verbatim}

The policy is a \emph{staircase threshold rule}: serve type 1 only when
$f_1 > f_0 + \theta(f_0)$ for a threshold $\theta$ that grows with $f_0$.
Beyond row $f_0 \approx d/2$, type 0 is never skipped regardless of $f_1$.

\subsection{Convergence of the Threshold to the Physical Deadline}

Define $\theta_0 = \min\{f_1 : \pi^*(0, f_1) = 1\}$ as the threshold at
$f_0 = 0$.  Table~\ref{tab:threshold} shows the convergence of
$\theta_0 / \text{tick\_rate}$ (in physical time units) toward $T_\text{real}$.

\begin{table}[h]
\centering
\caption{Threshold at $f_0 = 0$ in ticks and in physical time.}
\label{tab:threshold}
\begin{tabular}{rrrrr}
\toprule
tick\_rate & due$_\text{ticks}$ & $M$ & $\theta_0$ (ticks) & $\theta_0 /$ tick\_rate (physical) \\
\midrule
1 &  3 & 15 &  5 & 5.00 \\
2 &  6 & 25 &  7 & 3.50 \\
4 & 12 & 39 & 13 & 3.25 \\
8 & 24 & 71 & 24 & 3.00 \\
\bottomrule
\end{tabular}
\end{table}

The continuous-time limit of the threshold is $T_\text{real} = 3.0$ physical
time units: RVI serves type 1 only when type 1 has been waiting approximately
one full deadline duration longer than type 0.  This is the cost-rate trade-off
point: at that age difference, the per-unit-time urgency advantage of type 1
(accumulated waiting) compensates for the 3$\times$ cost-rate disadvantage.

\subsection{FIL-Projected Policy in the True k=4 System}
\label{sec:sectionA}

RVI is always solved on the FIL-only ($k=1$) projection, yet the policy is
evaluated in a richer simulator that tracks up to $k_\text{true} = 4$
positions per type.

\begin{table}[h]
\centering
\caption{FIFO vs.\ RVI (QL reward, tick\_rate = 1).}
\label{tab:sectionA}
\begin{tabular}{lcccc}
\toprule
System & FIFO mean & RVI mean & RVI/FIFO & Improvement \\
\midrule
$k = 1$ (FIL only) & 0.01123 & 0.01062 & 0.946 & 5.4\% \\
$k = 4$ (true system) & 0.01643 & 0.01477 & 0.899 & 10.1\% \\
\bottomrule
\end{tabular}
\end{table}

The absolute costs are higher in the $k=4$ system because the richer simulator
correctly accrues lateness for SIL/TIL/QIL jobs that the $k=1$ simulator
discards.  The RVI improvement is \emph{larger} at $k=4$ (10.1\% vs.\ 5.4\%):
the FIL projection transfers well because at $\rho = 0.25$ queue depths $>1$
are rare, while FIFO in both systems uses the same FIL-only information but
ignores cost rates entirely.

\subsection{Tick-Rate Granularity Sweep}

\begin{table}[h]
\centering
\caption{RVI vs.\ FIFO improvement across tick rates ($k=1$, QL reward).}
\label{tab:sectionC}
\begin{tabular}{rrrccr}
\toprule
tick\_rate & due$_\text{ticks}$ & $M$ & FIFO mean & RVI mean & Improvement \\
\midrule
1 &  3 & 15 & 0.01123 & 0.01062 &  5.4\% \\
2 &  6 & 25 & 0.01967 & 0.01809 &  8.0\% \\
4 & 12 & 39 & 0.03381 & 0.02944 & \textbf{12.9\%} \\
8 & 24 & 71 & 0.06511 & 0.06033 &  7.3\% \\
\bottomrule
\end{tabular}
\end{table}

Improvement increases monotonically from tick\_rate $= 1$ to $4$ (5.4\%
$\to$ 8.0\% $\to$ 12.9\%) and then regresses at tick\_rate $= 8$ (7.3\%).
Finer resolution gives RVI more decision points and a sharper threshold
(Table~\ref{tab:threshold} shows the physical-time threshold is coarser
at lower resolution).  The regression at tick\_rate $= 8$ is attributed to
the much larger state space ($M=71$, $\sim$111\,000 states): the g-stable
convergence criterion fires earlier relative to the Bellman-residual scale,
leaving $h$-values less well-converged and the near-deadline policy boundary
slightly imprecise.


% ============================================================
\section{Verification: Simulated Mean Cost vs.\ $g^*$}
\label{sec:verification}
% ============================================================

\subsection{Motivation}

Relative Value Iteration produces, upon convergence, two objects: the
average-cost estimate $g^*$ (a scalar) and the action map $\pi^*$ (a table of
optimal actions per state).  Both objects are derived from the same computation,
so internal consistency is guaranteed by construction.  A more meaningful
validation is to \emph{evaluate the action map independently}: simulate the
Markov chain under policy $\pi^*$, accumulate costs, and verify that the sample
mean matches $g^*$ to within Monte Carlo error.  This cross-check confirms that
(i) the Bellman iteration converged to the correct fixed point, (ii) the
action map was stored and queried correctly, and (iii) the simulation harness
itself is bug-free.

\subsection{Simulation Methodology (\texttt{EvaluatePolicyRaw})}

The function \texttt{EvaluatePolicyRaw} runs $\mathcal{T} = 300$ independent
trajectory replications, each of $N = 150\,000$ uniformised MDP steps, with a
warm-up of $W = 15\,000$ steps discarded at the start of each trajectory to
eliminate transient effects.

At each step the simulator proceeds as follows:
\begin{enumerate}
  \item If $\kappa = \textsc{AwaitAction}$: query $\pi^*(s)$ from the action map
        (using the encoded key; fallback: action 1 if key absent) and transition
        to the successor state.
  \item If $\kappa = \textsc{AwaitEvent}$: draw the next event type from the
        uniformised event distribution $(\lambda_0, \lambda_1, \gamma,
        \mu b_{k,j}, \ldots)$, advance the state, and accumulate cost $c(s)$.
\end{enumerate}

Three disjoint step counters are maintained over the post-warm-up horizon:
\begin{itemize}
  \item $N_\text{act}$: steps in \textsc{AwaitAction} (assignment decisions);
  \item $N_\text{ev}$: steps in \textsc{AwaitEvent} that draw a real event
        (arrival or service completion);
  \item $N_\text{tick}$: steps in \textsc{AwaitEvent} that draw a tick event
        (FIL refresh).
\end{itemize}
The total post-warm-up steps are $N_\text{total} = N_\text{act} + N_\text{ev} +
N_\text{tick}$.

\subsection{Two Notions of Average Cost}

\paragraph{Cost per RVI step.}
The quantity $g^*$ is defined as the average cost per step in the uniformised
chain.  Since cost is only incurred at \textsc{AwaitEvent} states, and since
action steps and tick steps carry zero immediate cost, the simulated analogue is:
\begin{equation}
  g^*_\text{sim}
    = \frac{\text{total accumulated cost}}{N_\text{total}}.
  \label{eq:g_per_step}
\end{equation}
This should agree with the Bellman $g^*$ to within Monte Carlo standard error.

\paragraph{Cost per real event.}
The policy comparer used in Section~\ref{sec:sectionA} and Table~\ref{tab:sectionC}
evaluates policies on a per-\emph{event} basis, where an event is an arrival or
service completion (not a tick and not an action step):
\begin{equation}
  g_\text{event}
    = \frac{\text{total accumulated cost}}{N_\text{ev}}.
  \label{eq:g_per_event}
\end{equation}

The two quantities are related by the \emph{step multiplier}
\begin{equation}
  \Lambda_\text{eff}
    = \frac{N_\text{total}}{N_\text{ev}}
    = \frac{N_\text{act} + N_\text{ev} + N_\text{tick}}{N_\text{ev}},
  \qquad
  g_\text{event} = \Lambda_\text{eff} \cdot g^*_\text{sim}.
  \label{eq:multiplier}
\end{equation}

\begin{remark}
  The step multiplier $\Lambda_\text{eff}$ can be interpreted as the ratio of
  the effective uniformisation rate $E[\Lambda]$ to the rate of real events.
  For tick\_rate $= 1$ and $\rho = 0.25$:
  \[
    E[\Lambda]
      = \lambda_0 + \lambda_1 + \gamma + \mu \cdot E[\text{busy servers}]
      = 0.1 + 0.1 + 1 + 0.4 \times 0.5
      = 1.4,
  \]
  \[
    E[\text{real event rate}]
      = \lambda_0 + \lambda_1 + \mu \cdot E[\text{busy}]
      = 0.2 + 0.2 = 0.4,
  \]
  giving $\Lambda_\text{eff} \approx 1.4 / 0.4 = 3.5$.  At tick\_rate $= 4$
  the tick contribution $\gamma = 4$ dominates and $\Lambda_\text{eff} \approx
  (0.4 + 4 + 0.2) / 0.4 \approx 11.5$.
\end{remark}

\subsection{Expected Outcome}

The simulation output is printed alongside the Bellman $g^*$ after convergence.
The key checks are:
\begin{enumerate}
  \item $g^*_\text{sim}$ should lie within $\approx 2\%$ of the Bellman $g^*$
        (Monte Carlo standard error decreases as $1/\sqrt{\mathcal{T} N}$).
  \item $g_\text{event}$ for the RVI policy should reproduce the ``RVI mean''
        entries of Table~\ref{tab:sectionA} and Table~\ref{tab:sectionC}.
  \item $g_\text{event}$ for the FIFO policy should reproduce the
        ``FIFO mean'' entries of the same tables.
  \item The step multiplier $\Lambda_\text{eff}$ printed from the step counter
        breakdown should match the theoretical formula~\eqref{eq:multiplier}.
\end{enumerate}

\begin{table}[h]
\centering
\caption{Structure of the verification output (QL reward, tick\_rate~$=1$,
         $k=1$, base configuration).}
\label{tab:verification}
\begin{tabular}{lll}
\toprule
Quantity & Symbol & Expected value \\
\midrule
Bellman convergence value & $g^*$ & (from RVI) \\
Simulated per-step mean, RVI & $g^*_\text{sim}$ & $\approx g^*$ ($<2\%$ error) \\
Simulated per-step mean, FIFO & --- & $> g^*$ (suboptimal) \\
Per-event mean, RVI & $g_\text{event}^\text{RVI}$ & matches Table~\ref{tab:sectionA} \\
Per-event mean, FIFO & $g_\text{event}^\text{FIFO}$ & matches Table~\ref{tab:sectionA} \\
Step multiplier & $\Lambda_\text{eff}$ & $\approx 3.5$ (tick\_rate $=1$) \\
\bottomrule
\end{tabular}
\end{table}

Agreement across all five checks closes the loop between the RVI Bellman
computation, the stored action map, the simulation harness, and the policy
comparer --- providing end-to-end confidence in the numerical implementation.


% ============================================================
\section{From Discrete Ticks to Continuous Time}
\label{sec:continuous}
% ============================================================

\subsection{Ticks Are Not Additional Decision Occasions}

A natural concern when interpreting a uniformised discrete-time MDP is whether
the discretisation introduces artificial decision opportunities that do not exist
in the underlying continuous-time system.  In a continuous-time queueing system
the scheduler is consulted only at \emph{events} --- arrivals and service
completions; between events the system evolves without intervention.

In the formulation studied here, tick events advance all FIL counters by one
unit and return the state to \textsc{AwaitEvent} \emph{without ever transitioning
to \textsc{AwaitAction}}.  The scheduler is therefore \emph{not consulted} at
tick events.

\begin{observation}
  The number of scheduling decisions equals the number of service completions
  plus the number of arrivals that find an idle server.  This is identical to
  the decision frequency in a continuous-time model.  Tick events create no
  additional scheduling decisions.
\end{observation}

The tick rate $\gamma$ therefore plays a purely \emph{informational} role: it
controls how finely the FIL values are measured, not how often the policy is
applied.

\subsection{The Real Gap: FIL Resolution}

The FIL value $f_n \in \{0, 1, 2, \ldots, M\}$ is an integer number of ticks.
In a continuous-time system the sojourn time of the FIL job would be a
non-negative real number $\tau_n$.  The two are related by
$\tau_n \approx f_n / \gamma$, with discretisation error $|\tau_n - f_n/\gamma|
\leq 1/(2\gamma)$.

Higher tick rates reduce this error:
\begin{itemize}
  \item At tick\_rate $= 1$: FIL is measured to the nearest 1.0 physical
        time unit.
  \item At tick\_rate $= 4$: FIL is measured to the nearest 0.25 physical
        time units.
  \item At tick\_rate $= 8$: FIL is measured to the nearest 0.125 physical
        time units.
  \item As $\gamma \to \infty$: $f_n / \gamma \to \tau_n$ in the limit; the
        FIL counter approaches a continuous sojourn timer.
\end{itemize}

The practical effect is visible in Table~\ref{tab:threshold}: the optimal
service threshold for type 1 at $f_0 = 0$, measured in physical time units
($\theta_0 / \gamma$), converges monotonically to $T_\text{real} = 3.0$ from
above as $\gamma$ increases.  At tick\_rate $= 8$ the threshold has already
reached $3.00$ to two decimal places.  This suggests that the continuous-time
optimal policy has a threshold precisely at the physical deadline --- a
structural result that may be derivable analytically via a fluid-limit or
Gittins-index argument.

\subsection{Translating $g^*$ to Continuous Time}

The Bellman average cost $g^*$ is a \emph{cost per uniformised step}.  Its
continuous-time analogue is a \emph{cost per physical time unit}:
\begin{equation}
  g^*_\text{time}
    \;=\; g^* \cdot E[\Lambda],
  \label{eq:g_time}
\end{equation}
where
\begin{equation}
  E[\Lambda]
    \;=\; \sum_n \lambda_n \;+\; \gamma \;+\; \mu \cdot E\!\Bigl[\textstyle\sum_{k,j} b_{k,j}\Bigr]
  \label{eq:elambda}
\end{equation}
is the expected uniformisation rate.  For the parameters used here and
tick\_rate $= r$:
\[
  E[\Lambda]
    = 0.2 + r + 0.4 \times 0.5
    = r + 0.4.
\]

The cost-per-event output of the policy comparer relates to $g^*_\text{time}$
by dividing by the rate of real events:
\begin{equation}
  g_\text{event}
    = \frac{g^*_\text{time}}{\lambda_\text{total} + \mu \cdot E[\text{busy}]}
    = \frac{g^* \cdot (r + 0.4)}{0.2 + 0.2}
    = g^* \cdot \frac{r + 0.4}{0.4}
    = g^* \cdot \Lambda_\text{eff},
  \label{eq:g_event_cont}
\end{equation}
consistent with the step multiplier formula~\eqref{eq:multiplier}.  Note that
$g^*_\text{time}$ is \emph{tick-rate-dependent}: as $\gamma$ increases, the
system accumulates cost more rapidly per unit of physical time (because FIL
values grow faster), even though the same physical dynamics are being modelled.
Results should therefore always be compared at the same tick rate.

\subsection{The Koole Approximation}

In the discrete model, the FIL of type $n$ is updated only at tick events.
When a service completion promotes the second-in-line job to FIL, the promoted
job's FIL value is inherited from the queue's running FIL counter --- not from
the exact time at which that job arrived.  The difference between the true
sojourn time and the FIL-approximate sojourn time is at most $1/\gamma$
(one tick interval).

This discrepancy is the \emph{Koole approximation}: the FIL counter is a
shared, coarsely-updated proxy for each job's individual age.  The approximation
improves as $\gamma$ increases.

\paragraph{Impact on decision quality.}
The scheduling policy acts on the FIL value of the oldest waiting job.  If the
true sojourn time is $\tau_n$ but the observed FIL implies $\tau_n \approx f_n
/ \gamma \pm 1/\gamma$, the policy uses a slightly imprecise age signal.  Near
the decision threshold ($\tau_n \approx T_\text{real}$) this $\pm 1/\gamma$
error can cause incorrect decisions for jobs right at the boundary.  At higher
tick rates the boundary is resolved more sharply and the decision error rate
decreases.

\subsection{Towards a Continuous-Time Simulator}

A continuous-time simulator would eliminate the Koole approximation entirely.
Rather than a FIL counter, each job in the queue carries its exact arrival time
$t_n^\text{arr}$.  At each event (arrival or completion), the exact sojourn
time of the oldest waiting job is:
\[
  \tau_n(t) = t - \min_{j \in \text{queue}_n} t_j^\text{arr}.
\]
The policy is queried with the continuous values $(\tau_0, \tau_1)$; to use the
trained RVI policy, these are discretised as $f_n = \lfloor \tau_n \cdot \gamma
\rfloor$ and the action map is looked up as usual.

Such a simulator would allow:
\begin{enumerate}
  \item \textbf{Direct verification of threshold convergence.}  By varying
        $\gamma$ across 1, 2, 4, 8, and querying the same trained policy, one
        can confirm that the serving threshold converges in physical time.
  \item \textbf{Quantification of the Koole approximation error.}  By comparing
        performance under exact sojourn times (continuous simulator) vs.\
        tick-approximate sojourn times (discrete MDP), the cost of the
        approximation can be isolated.
  \item \textbf{Asymptotic policy evaluation.}  At very high $\gamma$ the
        continuous simulator and the discrete MDP should agree; any residual
        difference is attributable to the $k=1$ FIL projection, not to the
        tick approximation.
\end{enumerate}

\subsection{The Discrete System as an Intermediate}

The study uses three levels of approximation simultaneously:
\begin{enumerate}
  \item \textbf{FIL projection ($k=1$):} only the oldest waiting job per type
        is tracked; all others are invisible to the policy.
  \item \textbf{Koole approximation:} FIL ages are tick-approximate, not exact.
  \item \textbf{Truncation at $M$:} states with $f_n > M$ are collapsed to $M$.
\end{enumerate}

These are nested approximations, each removing a source of information about
the queue.  Table~\ref{tab:sectionA} shows that the $k=1$ projection transfers
well to the $k_\text{true} = 4$ system at $\rho = 0.25$, because the probability
of having more than one job waiting is $O(\rho^2) \approx 6\%$.  Both the Koole
approximation and the truncation become negligible at high $\gamma$ and large
$M$, respectively.

The continuous-time extension (event-driven simulator, exact sojourn times,
all queue positions tracked) is identified as a natural follow-up, to be
implemented as a separate executable.  It would provide a clean upper bound on
what the FIL-only RVI policy can achieve in the true continuous system.


% ============================================================
\section{Why Binary Reward Produces Noisy Heatmaps}
\label{sec:binary}
% ============================================================

This section provides an elaborate explanation of the noisiness of heatmaps
under the binary deadline reward \eqref{eq:bin_cost}.  The argument proceeds
in five layers: the shape of the $h$-value landscape, the action-value gap,
structural near-ties, convergence amplification, and the persistence of noise
even under exact enumeration.

\subsection{Layer 1: The Cost Gradient Shapes the $h$-Value Landscape}

The relative value function $h(s)$ accumulates the long-run discounted
difference in cost between the policy $\pi^*$ and the reference policy.  Its
\emph{gradient} with respect to FIL reflects how rapidly the expected future
cost changes as a job ages.

\paragraph{QL reward.}
The per-step cost $c_n \cdot \max(0, f_n - d_n)$ increases by $c_n$ for every
additional tick beyond the deadline.  Consequently, for $f_n > d_n$:
\[
  \frac{\partial h_\text{QL}}{\partial f_n}
    \;\approx\; \frac{c_n}{\mu - \lambda_n}
    \;>\; 0.
\]
The $h$-value landscape is strictly increasing and \emph{everywhere has positive
gradient} for overdue jobs.  Even a one-tick difference in FIL translates into
a meaningful difference in $h$.

\paragraph{Binary reward.}
The per-step cost $c_n \cdot \mathbf{1}\{f_n > d_n\}$ is constant for all
$f_n > d_n$ and zero for all $f_n \leq d_n$.  The gradient is:
\[
  \frac{\partial h_\text{bin}}{\partial f_n}
    \;\approx\;
  \begin{cases}
    0 & f_n \ll d_n
        \quad (\text{safe region: negligible miss probability}), \\
    \text{large spike} & f_n \approx d_n
        \quad (\text{deadline cliff}), \\
    c_n \cdot p_\text{steady} & f_n \gg d_n
        \quad (\text{overdue: small constant increment}).
  \end{cases}
\]
Away from the deadline the $h$-value landscape is \emph{nearly flat}.

\subsection{Layer 2: The Action-Value Gap}

At an \textsc{AwaitAction} state $s$ with $f_1 > f_0$, the two action values
are
\[
  Q(s, a) = \sum_{s'} P(s' \mid s, a)\, h(s'), \qquad a \in \{0,1\}.
\]
The \emph{action-value gap}
\[
  m(s) = Q(s, 0) - Q(s, 1)
\]
determines the policy: $m > 0$ means action 1 is optimal (assign FIFO-top
job), $m < 0$ means action 0 is optimal (skip FIFO-top job to serve the
other type), $m \approx 0$ is a near-tie.

\paragraph{QL reward (large gap).}
For $f_1 > d_1$ and $f_0 < d_0$, assigning service to type 1 immediately
removes its per-step cost of $c_1 \cdot (f_1 - d_1)$ per step, whereas
deferring type 1 incurs this cost for approximately $1/\mu$ additional steps.
The gap is approximately
\[
  m_\text{QL}(s)
    \;\approx\; c_1 \cdot \frac{f_1 - d_1}{\mu}
    \;-\; c_0 \cdot \frac{f_0 - d_0^+}{\mu},
\]
which is $O(c_n / \mu)$ and grows linearly with the overdue-ness of each type.
For the parameters used here ($c_1 = 1$, $\mu = 0.4$, overdue-ness $\approx
1$--$10$ ticks) the gap is in the range $[2.5, 25]$.

\paragraph{Binary reward (small or zero gap).}
The action gap under binary cost can be expressed as
\[
  m_\text{bin}(s)
    = c_1 \cdot \bigl[P(\text{type 1 misses} \mid a=0) - P(\text{type 1 misses} \mid a=1)\bigr]
    - c_0 \cdot \bigl[P(\text{type 0 misses} \mid a=0) - P(\text{type 0 misses} \mid a=1)\bigr].
\]
\begin{itemize}
  \item For $f_0, f_1 \ll d$ (both jobs are safe): $P(\text{miss}) \approx 0$
        for either action, so $m_\text{bin}(s) \approx 0$.
  \item For $f_1 \approx d$ (type 1 at the deadline): $P(\text{type 1 misses}
        \mid a=0)$ can be $O(0.1)$--$O(1)$, making the gap non-trivial but
        small relative to the total state space.
  \item For $f_1 \gg d$ (type 1 long overdue): both actions incur the binary
        cost with high probability, and the gap is again small.
\end{itemize}

\subsection{Layer 3: The Extent of the Near-Tie Region}

Under QL, the only near-ties are states near the \emph{boundary} of the optimal
threshold staircase.  These are a thin band of cells in the $(f_0, f_1)$ grid;
the interior is decisive.

Under binary, near-ties dominate:
\begin{itemize}
  \item \textbf{Safe region} ($f_0, f_1 < d$):
        both jobs are unlikely to miss the deadline under any reasonable policy.
        The gap is proportional to the (tiny) probability that either job
        reaches the deadline before being served.  This region covers roughly
        $d^2 = d_\text{ticks}^2$ cells in the grid.  For tick\_rate $= 4$ and
        $d = 12$, this is $144$ cells, all with near-zero gap.

  \item \textbf{Far-overdue region} ($f_0, f_1 \gg d$):
        both jobs have already missed the deadline.  The binary cost fires for
        both regardless of the action, and the gap is again small.

  \item \textbf{Decisive region} ($f_1 \approx d$, one or both types near the
        deadline cliff): the gap is meaningful, but this is a thin strip.
\end{itemize}

The fraction of cells with a decisive gap is thus $O(d / M)$, typically 20--30\%
of the grid.  The remaining 70--80\% of cells have near-zero gaps and ambiguous
policies.

\subsection{Layer 4: Convergence Amplification}

The \texttt{g\_stable} stopping criterion guarantees that
$\lvert g^{(t+1)} - g^{(t)} \rvert < \varepsilon_g = 10^{-10}$, but it
\emph{does not} guarantee that $h^{(t)}$ has converged to within $\varepsilon_g$
of the true $h^*$.  In fact, at termination the span of the residual
$h^{(t)} - h^*$ is typically $O(\varepsilon_\text{span})$, where
$\varepsilon_\text{span}$ is the span at termination---which may be much larger
than $\varepsilon_g$.

A near-tie is resolved \emph{incorrectly} whenever
\begin{equation}
  \lvert m(s) \rvert \;<\; 2 \cdot \varepsilon_\text{span}.
  \label{eq:neartie}
\end{equation}
Under QL, the typical action gap is $O(c_n / \mu) \approx 2$--$25$, far larger
than $\varepsilon_\text{span} \approx 10^{-3}$--$10^{-1}$ at \texttt{g\_stable}
termination.  Condition~\eqref{eq:neartie} is satisfied only at the true policy
boundary---a thin band.

Under binary, the typical action gap away from the deadline cliff is
$O(10^{-4})$--$O(10^{-2})$, \emph{comparable to or smaller than}
$\varepsilon_\text{span}$.  Condition~\eqref{eq:neartie} is satisfied for most
of the safe and far-overdue regions.  The argmin stored in \texttt{action\_map}
for these cells was determined not by the true optimal policy but by the sign of
a tiny numerical residual in $h^{(t)}$.

\paragraph{Scale mismatch.}
The problem is compounded by the scale of the $h$-values.  Under QL with
$\text{tick\_rate} = 4$ and $M = 39$, values of $h$ can reach
$c_0 \cdot (M - d) / (\mu - \lambda_0) \approx 3 \times 27 / 0.3 \approx 270$.
Under binary with the same parameters, $h$ is bounded by
$c_0 \cdot P(\text{miss}) \lesssim 3 \times 1 = 3$.
The \emph{relative} convergence error $\varepsilon_\text{span} / \|h\|_\infty$
is thus roughly 90 times larger under binary than under QL for the same
absolute convergence tolerance.

\subsection{Layer 5: Noise Persists Under Enumeration}

The simulation-based heatmap has two noise sources: (a) unvisited cells, and
(b) cells sampled under different busy configurations $B$.  The enumerated
heatmap eliminates both by directly constructing the canonical state.

However, the enumerated heatmap can only report what is stored in
\texttt{action\_map}.  If an \texttt{action\_map} entry was itself determined
by a near-tie in the imperfectly converged $h^{(t)}$, the enumerated cell
faithfully displays the \emph{wrong} argmin.  No amount of sampling improvement
fixes this.

Concretely, suppose at state $s$ (with $f_0 = 5$, $f_1 = 5$, safely below the
deadline) the true gap is $m^*(s) = +0.0003$ (slightly prefer to assign type 1)
but the residual in $h^{(t)}$ introduces an error $\delta = -0.0005$.  The
stored argmin is action 0 (skip), the wrong choice.  The enumerated heatmap
will display ``\texttt{.}'' at this cell, creating an isolated artefact in an
otherwise smooth region.

The result is a \emph{salt-and-pepper} pattern in the heatmap:
\begin{itemize}
  \item Large connected regions of type-0 (below the diagonal) are correct,
        since $m \gg 0$ there.
  \item Near the diagonal, isolated \texttt{.} and \texttt{1} cells flip
        unpredictably based on which direction the convergence residual
        happened to point.
  \item There is no clean staircase boundary; the boundary is a ragged,
        jagged band one to three cells wide.
\end{itemize}

\subsection{Summary: QL vs.\ Binary}

\begin{table}[h]
\centering
\caption{Comparison of QL and binary rewards for heatmap quality.}
\label{tab:compare}
\begin{tabular}{lcc}
\toprule
Property & QL reward & Binary reward \\
\midrule
Cost gradient w.r.t.\ FIL (overdue) & Linear, always $> 0$ & Flat; spike only at $d$ \\
Typical action gap $|m(s)|$ & $O(c_n / \mu) \sim 2$--$25$ & $O(10^{-3})$ away from cliff \\
Fraction of near-tie cells & $\sim$5--10\% (boundary only) & $\sim$70--80\% (safe + far-overdue) \\
Scale of $h$-values & $O(c_n \cdot M / \mu) \sim 100$ & $O(c_n \cdot P(\text{miss})) \sim 1$ \\
Relative convergence error & Small & Large \\
Heatmap under simulation & Clean with sampling gaps & Noisy everywhere \\
Heatmap under enumeration & Clean, artefact-free & Ragged boundary, isolated flips \\
\bottomrule
\end{tabular}
\end{table}

In summary, the QL reward produces clean heatmaps because its smooth,
everywhere-positive cost gradient generates large, decisive action-value gaps
that dwarf convergence residuals.  The binary reward produces noisy heatmaps
because its near-zero gradient in the safe and far-overdue regions creates
pervasive near-ties that convergence residuals --- unavoidable under the
\texttt{g\_stable} fallback --- resolve arbitrarily.  The noise is a structural
property of the binary cost landscape, not merely a simulation artefact.


% ============================================================
\section{Discussion}
\label{sec:discussion}
% ============================================================

\paragraph{The threshold as a structural result.}
The convergence of the type-1 serving threshold to $T_\text{real}$ in physical
time (Table~\ref{tab:threshold}) is a computationally derived structural result.
It suggests that the optimal policy for this system is approximated by a simple
parametric rule:
\[
  \pi^*(f_0, f_1)
  = \begin{cases}
      \text{serve type 0} & \text{if } f_0 \geq f_1, \\
      \text{serve type 1} & \text{if } f_1 > f_0 + d + \epsilon(f_0), \\
      \text{serve type 0} & \text{otherwise},
    \end{cases}
\]
where $\epsilon(f_0) \geq 0$ is a small correction term that grows with $f_0$.
Whether such a rule can be derived analytically (e.g.\ via Gittins indices or a
fluid-limit argument) is an open question.

\paragraph{Fixing binary-reward convergence.}
Two remedies exist for the binary case:
\begin{enumerate}
  \item \textbf{More iterations / span criterion.}  Running RVI until
        $\text{span} < \varepsilon_\text{span}$ rather than using the
        \texttt{g\_stable} fallback would guarantee sufficiently converged
        $h$-values.  The required number of iterations may be very large for
        binary cost because the information about cost must propagate
        ``backward'' across the many near-zero-gradient states between the
        deadline and the current FIL.

  \item \textbf{Regularisation.}  Adding a small QL-like component to the
        binary cost (e.g.\ $c_n \cdot \max(0, f_n - d_n) \cdot \varepsilon_\text{reg}$)
        ensures a positive gradient everywhere, restoring large action gaps.
        As $\varepsilon_\text{reg} \to 0$ the solution approaches the true
        binary optimum.
\end{enumerate}


% ============================================================
\section{Conclusion}
\label{sec:conclusion}
% ============================================================

We have presented an enumeration-based heatmap that queries the RVI action map
directly for a canonical state at each $(f_0, f_1)$ cell, eliminating all
simulation artefacts.  Under the queue-lateness reward the heatmap exposes a
clean three-region staircase policy: always serve the cheaper type when it is
competitive on age, and only defer to the costlier type when its age excess
exceeds approximately one physical deadline $T_\text{real}$.  This threshold
converges to $T_\text{real} = 3.0$ as the tick rate increases from 1 to 8.

The Bellman average-cost value $g^*$ was validated end-to-end via an
independent simulation (Section~\ref{sec:verification}): the
\texttt{EvaluatePolicyRaw} harness simulates the policy over 300 replications
of 150\,000 steps and confirms that the per-step sample mean matches $g^*$ to
within Monte Carlo error.  The step multiplier $\Lambda_\text{eff}$ derived
from the step-counter breakdown bridges the per-RVI-step and per-real-event
cost denominators, reconciling the Bellman value with the policy-comparer
output.

The discrete uniformised chain relates to the continuous-time system in two
distinct ways (Section~\ref{sec:continuous}): tick events create \emph{no}
additional scheduling decisions (the decision frequency is identical to
continuous time), while the tick rate controls the \emph{resolution} of FIL
measurement.  The threshold convergence in physical time
(Table~\ref{tab:threshold}) is therefore a statement about the continuous-time
optimal policy, not an artefact of discretisation.  A continuous-time
event-driven simulator tracking exact sojourn times would eliminate the Koole
approximation and provide a clean benchmark; this is identified as future work.

The central technical contribution of this note is a precise explanation of
why the binary reward produces noisy heatmaps.  The binary cost's near-zero
gradient away from the deadline cliff creates action-value near-ties across
70--80\% of the state space.  These near-ties are resolved by the sign of the
RVI convergence residual --- which the \texttt{g\_stable} stopping criterion
does not control tightly enough --- producing a salt-and-pepper artefact that
is a property of the stored action map itself, not of the heatmap rendering
method.  Switching to the QL reward, or running RVI to a tighter span
tolerance, eliminates the noise.

\end{document}
